% GENERATED FILE. Edit canonical lessons, then run npm run generate:books.
% canonical-source-hash: fnv1a64:10eb405f40c6c035
% canonical-lessons: KA-C138-rupayi, KA-C138-batali, KA-C138-jebu, KA-C138-saks, KA-C138-kotu

\chapter{Rupee, Bottle, Pocket, Socks, Coat}
\label{ch:ka-rupee-bottle-pocket-socks-coat}
\hlchaptermodality{138}

\begin{chapteropening}
\textbf{By the end of this chapter:} \emph{I can say a rupee, a bottle, a pocket, socks and a coat.}

Complete the last lesson of chapter 138: I can say a rupee, a bottle, a pocket, socks and a coat.
\end{chapteropening}

\section[\texorpdfstring{\kn{ರೂಪಾಯಿ} (rūpāyi)}{ರೂಪಾಯಿ (rūpāyi)}]{\kn{ರೂಪಾಯಿ} (rūpāyi) — a rupee}

\label{lesson:KA-C138-rupayi}



\begin{quote}
Before the new one: say the Kannada for a television, then the Kannada for a computer.

\emph{Read it:} the digits aloud — \textbf{\kn{೨}}, \textbf{\kn{೪}}, \textbf{\kn{೫}}, \textbf{\kn{೬}}, \textbf{\kn{೭}}, \textbf{\kn{೮}}, \textbf{\kn{೯}}, \textbf{\kn{೦}}
\end{quote}

\subsection*{You'll want to know: \kn{ರೂಪಾಯಿ}}
\textbf{\kn{ರೂಪಾಯಿ}} — \emph{rūpāyi} — \textquotedblleft{}a rupee\textquotedblright{}.

\begin{grammarlens}[title={What you've built}]
One more word for this chapter.
\end{grammarlens}

\subsection*{Your turn}
\begin{itemize}
\raggedright
  \item \emph{Say it:} \emph{rūpāyi}
  \item \emph{Say it:} \emph{rūpāyi}, once more
  \item \emph{Say it:} say \emph{rūpāyi}
  \item \emph{Recall:} say the Kannada for a television, then the Kannada for a computer, then say \emph{rūpāyi} again
\end{itemize}

\begin{tcolorbox}[breakable,colback=teal!4,colframe=teal!35!black,title={Before you move on}]
What does \kn{ರೂಪಾಯಿ} mean? (\textquotedblleft{}A rupee\textquotedblright{}.) Say it once more.
\end{tcolorbox}
\section[\texorpdfstring{\kn{ಬಾಟಲಿ} (bāṭali)}{ಬಾಟಲಿ (bāṭali)}]{\kn{ಬಾಟಲಿ} (bāṭali) — a bottle}

\label{lesson:KA-C138-batali}



\begin{quote}
Before the new one: say the Kannada for a computer, then the Kannada for a rupee.
\end{quote}

\subsection*{You'll want to know: \kn{ಬಾಟಲಿ}}
\textbf{\kn{ಬಾಟಲಿ}} — \emph{bāṭali} — \textquotedblleft{}a bottle\textquotedblright{}.

\begin{grammarlens}[title={What you've built}]
One more word for this chapter.
\end{grammarlens}

\subsection*{Your turn}
\begin{itemize}
\raggedright
  \item \emph{Say it:} \emph{bāṭali}
  \item \emph{Say it:} \emph{bāṭali}, once more
  \item \emph{Say it:} say \emph{bāṭali}
  \item \emph{Recall:} say the Kannada for a computer, then the Kannada for a rupee, then say \emph{bāṭali} again
\end{itemize}

\begin{tcolorbox}[breakable,colback=teal!4,colframe=teal!35!black,title={Before you move on}]
What does \kn{ಬಾಟಲಿ} mean? (\textquotedblleft{}A bottle\textquotedblright{}.) Say it once more.
\end{tcolorbox}
\section[\texorpdfstring{\kn{ಜೇಬು} (jēbu)}{ಜೇಬು (jēbu)}]{\kn{ಜೇಬು} (jēbu) — a pocket}

\label{lesson:KA-C138-jebu}



\begin{quote}
Before the new one: say the Kannada for a rupee, then the Kannada for a bottle.
\end{quote}

\subsection*{You'll want to know: \kn{ಜೇಬು}}
\textbf{\kn{ಜೇಬು}} — \emph{jēbu} — \textquotedblleft{}a pocket\textquotedblright{}.

\begin{grammarlens}[title={What you've built}]
One more word for this chapter.
\end{grammarlens}

\subsection*{Your turn}
\begin{itemize}
\raggedright
  \item \emph{Say it:} \emph{jēbu}
  \item \emph{Say it:} \emph{jēbu}, once more
  \item \emph{Say it:} say \emph{jēbu}
  \item \emph{Recall:} say the Kannada for a rupee, then the Kannada for a bottle, then say \emph{jēbu} again
\end{itemize}

\begin{tcolorbox}[breakable,colback=teal!4,colframe=teal!35!black,title={Before you move on}]
What does \kn{ಜೇಬು} mean? (\textquotedblleft{}A pocket\textquotedblright{}.) Say it once more.
\end{tcolorbox}
\section[\texorpdfstring{\kn{ಸಾಕ್ಸ್} (sāks)}{ಸಾಕ್ಸ್ (sāks)}]{\kn{ಸಾಕ್ಸ್} (sāks) — socks}

\label{lesson:KA-C138-saks}



\begin{quote}
Before the new one: say the Kannada for a bottle, then the Kannada for a pocket.
\end{quote}

\subsection*{You'll want to know: \kn{ಸಾಕ್ಸ್}}
\textbf{\kn{ಸಾಕ್ಸ್}} — \emph{sāks} — \textquotedblleft{}socks\textquotedblright{}.

\begin{grammarlens}[title={What you've built}]
One more word for this chapter.
\end{grammarlens}

\subsection*{Your turn}
\begin{itemize}
\raggedright
  \item \emph{Say it:} \emph{sāks}
  \item \emph{Say it:} \emph{sāks}, once more
  \item \emph{Say it:} say \emph{sāks}
  \item \emph{Recall:} say the Kannada for a bottle, then the Kannada for a pocket, then say \emph{sāks} again
\end{itemize}

\begin{tcolorbox}[breakable,colback=teal!4,colframe=teal!35!black,title={Before you move on}]
What does \kn{ಸಾಕ್ಸ್} mean? (\textquotedblleft{}Socks\textquotedblright{}.) Say it once more.
\end{tcolorbox}
\section[\texorpdfstring{\kn{ಕೋಟು} (kōṭu)}{ಕೋಟು (kōṭu)}]{\kn{ಕೋಟು} (kōṭu) — a coat}

\label{lesson:KA-C138-kotu}



\begin{quote}
Before the new one: say the Kannada for a pocket, then the Kannada for socks.
\end{quote}

\subsection*{You'll want to know: \kn{ಕೋಟು}}
\textbf{\kn{ಕೋಟು}} — \emph{kōṭu} — \textquotedblleft{}a coat\textquotedblright{}.

\begin{grammarlens}[title={What you've built}]
One more word for this chapter.
\end{grammarlens}

\subsection*{Your turn}
\begin{itemize}
\raggedright
  \item \emph{Say it:} \emph{kōṭu}
  \item \emph{Say it:} \emph{kōṭu}, once more
  \item \emph{Say it:} say \emph{kōṭu}
  \item \emph{Recall:} say the Kannada for a pocket, then the Kannada for socks, then say \emph{kōṭu} again
\end{itemize}

\begin{tcolorbox}[breakable,colback=teal!4,colframe=teal!35!black,title={Before you move on}]
What does \kn{ಕೋಟು} mean? (\textquotedblleft{}A coat\textquotedblright{}.) Say it once more.
\end{tcolorbox}
